\documentclass[11pt]{scrartcl}
\usepackage{ebb}



\title {Problemas Diarios}
\subtitle {Ponte a entrenar ps}

\author{Emmanuel Buenrostro}

\begin{document}
\maketitle
\tableofcontents
\problemdiff{2}
\problemtags{Combi, Invarianza, Coloraciones}
\begin{problema} [OMM 2017/1]
    
    En un tablero de ajedrez de $2017 \times 2017$, se han colocado en la primera columna $2017$
caballos de ajedrez, uno en cada casilla de la columna. Una tirada consiste en elegir dos
caballos distintos y de manera simulátnea moverlos como se mueven los caballos de ajedrez.
Encuentra todos los posibles valores enteros de $k$ con $1 \leq k \leq 2017$, para los cuales es
posible llegar a través de varias tiradas, a que todos los caballos están en la columna $k$, uno
en cada casilla.
Nota. Un caballo se mueve de una casilla $X$ a otra $Y$, solamente si $X$ y $Y$ son las esquinas
opuestas de un rectángulo de $3 \times 2$ o de $2 \times 3$.
    
\end{problema}
\problemdiff{4}
\problemtags{Combi, Invarianza, Coloraciones}
\begin{problema} [OMM 2017/1]
    
    En un tablero de ajedrez de $2017 \times 2017$, se han colocado en la primera columna $2017$
caballos de ajedrez, uno en cada casilla de la columna. Una tirada consiste en elegir dos
caballos distintos y de manera simulátnea moverlos como se mueven los caballos de ajedrez.
Encuentra todos los posibles valores enteros de $k$ con $1 \leq k \leq 2017$, para los cuales es
posible llegar a través de varias tiradas, a que todos los caballos están en la columna $k$, uno
en cada casilla.
Nota. Un caballo se mueve de una casilla $X$ a otra $Y$, solamente si $X$ y $Y$ son las esquinas
opuestas de un rectángulo de $3 \times 2$ o de $2 \times 3$.
    
\end{problema}

\problemdiff{7}
\problemtags{Combi, Invarianza, Coloraciones}
\begin{problema} [OMM 2017/1]
    
    En un tablero de ajedrez de $2017 \times 2017$, se han colocado en la primera columna $2017$
caballos de ajedrez, uno en cada casilla de la columna. Una tirada consiste en elegir dos
caballos distintos y de manera simulátnea moverlos como se mueven los caballos de ajedrez.
Encuentra todos los posibles valores enteros de $k$ con $1 \leq k \leq 2017$, para los cuales es
posible llegar a través de varias tiradas, a que todos los caballos están en la columna $k$, uno
en cada casilla.
Nota. Un caballo se mueve de una casilla $X$ a otra $Y$, solamente si $X$ y $Y$ son las esquinas
opuestas de un rectángulo de $3 \times 2$ o de $2 \times 3$.
    
\end{problema}

\problemdiff{9}
\problemtags{Combi, Invarianza, Coloraciones}
\begin{problema} [OMM 2017/1]
    
    En un tablero de ajedrez de $2017 \times 2017$, se han colocado en la primera columna $2017$
caballos de ajedrez, uno en cada casilla de la columna. Una tirada consiste en elegir dos
caballos distintos y de manera simulátnea moverlos como se mueven los caballos de ajedrez.
Encuentra todos los posibles valores enteros de $k$ con $1 \leq k \leq 2017$, para los cuales es
posible llegar a través de varias tiradas, a que todos los caballos están en la columna $k$, uno
en cada casilla.
Nota. Un caballo se mueve de una casilla $X$ a otra $Y$, solamente si $X$ y $Y$ son las esquinas
opuestas de un rectángulo de $3 \times 2$ o de $2 \times 3$.
    
\end{problema}

\printindex[topic]

\end{document}
